\documentclass[12pt]{article}
\usepackage[T1]{fontenc}
\usepackage[utf8]{inputenc}
\usepackage{lmodern}
\usepackage{textcomp}
\usepackage{enumitem}
\usepackage{geometry}
\usepackage{graphicx}
\usepackage{amsmath, amssymb}
\geometry{margin=1in}
\begin{document}
\begin{center}
Economics - K-12 Level Assessment 12 \
Total Marks: 40 \
Answer all questions.
\end{center}

\section*{Multiple Choice (10 marks)}
\begin{enumerate}[label=\arabic*.]
\item If consumers are advised that multigrained bread will substantially lessen the risk of cancer, which of the following will happen in the market for multigrained bread?
\begin{enumerate}[label=(\alph*)]
    \item The demand curve will shift to the left, decreasing the price of multigrained bread.
    \item The supply curve will shift to the left, increasing the price of multigrained bread.
    \item The demand curve will shift to the right, increasing the price of multigrained bread.
    \item The supply curve will shift to the right, decreasing the price of multigrained bread.
\end{enumerate}
\item The competitive market for gasoline, a normal good, is currently in a state of equilibrium. Which of the following would most likely increase the price of gasoline?
\begin{enumerate}[label=(\alph*)]
    \item Household income falls.
    \item Technology used to produce gasoline improves.
    \item The price of subway tickets and other public transportation falls.
    \item The price of crude oil, a raw material for gasoline, rises.
\end{enumerate}
\item Suppose the President plans to cut taxes for consumers and also plans to increase spending on the military. How does this affect real GDP and the price level?
\begin{enumerate}[label=(\alph*)]
    \item GDP increases and the price level decreases.
    \item GDP decreases and the price level increases.
    \item GDP stays the same and the price level increases.
    \item GDP increases and the price level increases.
\end{enumerate}
\item In the equation of exchange if V and Q are constant then
\begin{enumerate}[label=(\alph*)]
    \item changes in the price level must be proportional to changes in the money supply.
    \item changes in the money supply have no effect on the price level.
    \item changes in the price level have no effect on the money supply.
    \item the equation is invalid.
\end{enumerate}
\item The standard of living is measured by
\begin{enumerate}[label=(\alph*)]
    \item GDP.
    \item GDP per capita.
    \item Real GDP per capita.
    \item Actual GI)P per capita.
\end{enumerate}
\end{enumerate}

\section*{True / False (10 marks)}
\textit{Answer True or False}
\begin{enumerate}[label=\arabic*.]
\setcounter{enumi}{5}
\item True or False: If the price of aluminum falls, the supply of bicycles is likely to increase, assuming other factors remain constant.
\item True or False: According to Keynesian analysis, a decrease in the money supply lowers both the price level and output in the economy.
\item True or False: If global warming eliminates snow, ski instructors would face structural unemployment due to a permanent change in the job market.
\item True or False: Business cycles occur infrequently and do not follow a predictable pattern in capitalist economies.
\item True or False: Budget surpluses and higher discount rates are effective strategies to combat economic contraction during a business cycle.
\end{enumerate}

\section*{Long Form Response (20 marks)}
\textit{Answer each question with a clear and complete explanation.}
\begin{enumerate}[label=\arabic*.]
\setcounter{enumi}{10}
\item Discuss how a firm's ability to increase output more than proportionately to its inputs, exemplified by tripling inputs and quadrupling outputs, illustrates the concept of economies of scale.
\item Discuss how lowering personal income taxes could affect the aggregate demand in an economy. What are the potential consequences for consumers and overall economic activity?
\end{enumerate}

\vfill
\noindent\textit{End of Paper}
\end{document}